\textbf{Proof.}
By \(\mathrm{lem:source\text{-}only\text{-}frontier\text{-}lower}\), for the fixed model parameters there are \(c_0>0\) and \(n_0\) such that, for every \(n\ge n_0\) and every \(m\),
\[
\mathfrak R_{n,m}(a,s,t)\ge c_0\rho_n(a,s).
\tag{1}
\]
The lower-bound subexperiments in that lemma use \(f_T=1\). Since the declared target-density constants satisfy \(c_{T,-}\le1\le c_{T,+}\), those subexperiments remain in the model when either target-density bound is attained. Thus (1) uses no strict slack in the target-density bounds.

By \(\mathrm{ass:sample\text{-}balance}\), there are constants \(c_\kappa,C_\kappa>0\) and \(n_1\) such that
\[
c_\kappa n\le m\le C_\kappa n,
\qquad n\ge n_1.
\tag{2}
\]
In particular,
\[
m^{-1/2}\le c_\kappa^{-1/2}n^{-1/2}.
\tag{3}
\]
The regular branch satisfies \(\rho_n(a,s)=n^{-1/2}\). In the critical branch,
\[
\rho_n(1,s)=\sqrt{\log n/n}\ge n^{-1/2}
\tag{4}
\]
eventually. In the nonregular branch, \(s>0\) and \(a>1\) imply
\[
\frac{s+1}{2s+a+1}<\frac12,
\qquad
\rho_n(a,s)=n^{-(s+1)/(2s+a+1)}\ge n^{-1/2}.
\tag{5}
\]
It follows from (3)--(5) that, for some finite \(C_1\),
\[
m^{-1/2}+\rho_n(a,s)\le C_1\rho_n(a,s)
\tag{6}
\]
eventually. Combining (1) and (6) gives
\[
\mathfrak R_{n,m}(a,s,t)
\ge \frac{c_0}{C_1}\{m^{-1/2}+\rho_n(a,s)\}.
\tag{7}
\]
Substitution of the three branches from \(\mathrm{def:frontier\text{-}rate}\) proves the displayed minimax lower phase diagram.

The corrected \(\mathrm{lem:oracle\text{-}truncation\text{-}envelope}\) applies to every law in the same model class. It gives both the actual-bias inequality
\[
|b_{h_n}(\mu;P)|\le b_{h_n}^o(P)\le Ch_n^{s+1}
\tag{8}
\]
and the deterministic oracle comparison
\[
m^{-1/2}+h_n^{s+1}
+n^{-1/2}\|\omega_{h_n}^o\|_{L^2(P_S)}
\asymp m^{-1/2}+\rho_n(a,s).
\tag{9}
\]
That lemma also proves that the allowance in (8) is class-sharp for \(a>0\), whereas at \(a=0\), with \(\alpha=\min\{t,1\}\),
\[
\sup_P b_{h_n}^o(P)\le C h_n^{s+1+\alpha}.
\tag{10}
\]
Even at this endpoint, \(h_n^{s+1}=n^{-1/2}\), so retaining \(h_n^{s+1}\) in (9) does not enlarge the order of the regular oracle benchmark.

Finally, \(\omega_{h_n}^o\) depends on population integrals and the true density ratio. Thus (9) is a boundary-corrected oracle linear-weight benchmark, not an estimator measurable with respect to \(\mathcal D_{n,m}\). It cannot be used as a minimax upper bound or as an attainment claim. Equations (7)--(10) prove the theorem. \(\square\)